\chapter{검증: 기준 모델과 대조한다}\label{ch:verify}

이 장은 회로가 옳다는 것을 어떻게 확인하는지를 다룬다. 내용은 세 가지다. 같은 HDL로 쓰였어도 역할이 다른
두 종류의 코드, 기준 모델과 대조하는 검증 방식, 그리고 서로 속도가 다른 블록을 잇는 약속(핸드셰이크)이다.

\section{시뮬레이션용 모델과 합성용 설계}\label{sec:sim-vs-synth}

\Why HDL 코드는 두 가지 방식으로 쓰인다. \term{시뮬레이터}{simulator}는 코드를 컴퓨터 위에서 흉내 내어
신호가 클럭마다 어떻게 변하는지 계산한다. 합성 도구는 같은 코드를 실제 부품으로 바꾼다. 시뮬레이터는
언어가 허용하는 것은 무엇이든 계산해 주지만, 합성 도구는 부품으로 바꿀 수 있는 것만 받아들이고, 받아들인
것도 터무니없이 크거나 느린 회로가 될 수 있다. 그래서 ``시뮬레이션에서 옳게 동작하는 코드''와 ``칩에
올릴 수 있는 코드'' 사이에는 간격이 있다.

이 간격을 실수가 아니라 도구로 쓸 수 있다. 먼저 크기와 속도를 생각하지 않고 \emph{무엇을 해야 하는지}만
분명하게 적은 코드를 만든다. 이것이 시뮬레이션용 기준 모델이다. 그다음 같은 동작을 하되 칩에 맞게 구조를
바꾼 합성용 설계를 만들고, 둘의 출력이 같은지 비교한다. 동작의 정의와 구현의 최적화를 분리하는 것이다.

\General FIFO(먼저 넣은 것이 먼저 나오는 대기열)를 예로 들자. 시뮬레이션용 모델은 몇 줄이면 된다. 언어가
제공하는 가변 길이 목록에 넣고 빼면 된다. 이 코드는 합성할 수 없다. 길이가 변하는 목록에 해당하는 부품이
없기 때문이다. 합성용 FIFO는 고정된 크기의 메모리와 쓰기 위치, 읽기 위치를 가리키는 두 카운터로 만든다
(\ref{sec:handshake}절). 두 코드는 생긴 것이 전혀 다르지만, 같은 순서로 넣고 빼면 같은 값이 나와야 한다.

표~\ref{tab:g-synth}는 시뮬레이션에서는 문제없지만 합성에서는 조심해야 하는 것들이다.

\begin{table}[ht]
\centering
\caption{시뮬레이션과 합성에서 다르게 다뤄지는 코드.}
\label{tab:g-synth}
\small
\begin{tabularx}{\linewidth}{>{\raggedright\arraybackslash}p{42mm}>{\raggedright\arraybackslash}X>{\raggedright\arraybackslash}X}
\toprule
코드 & 시뮬레이션 & 합성 \\
\midrule
시간 지연(\code{\#10}), 화면 출력 & 그대로 수행 & 무시되거나 오류 \\
가변 길이 목록, 동적 할당 & 그대로 수행 & 불가능 \\
상수가 아닌 수로 나누기, 나머지 & 그대로 계산 & 큰 나눗셈 회로 \\
변수로 큰 배열의 원소 고르기 & 배열 접근 한 번 & 원소 수만큼 큰 선택 회로 \\
큰 반복문 & 금방 끝남 & 본체가 반복 횟수만큼 복사됨 \\
\bottomrule
\end{tabularx}
\end{table}

\Ours 이 프로젝트에는 모델이 세 층으로 있다(그림~\ref{fig:g-layers}).

\begin{figure}[ht]
\centering
\begin{tikzpicture}[x=1cm,y=1cm]
  \node[blkg,minimum width=36mm,minimum height=14mm] (rust) at (0,0) {Rust 코어\\{\footnotesize 무엇을 계산하는가}};
  \node[blkb,minimum width=36mm,minimum height=14mm] (ref) at (5,0) {기준 RTL\\{\footnotesize 클럭 단위로 어떻게}};
  \node[blko,minimum width=36mm,minimum height=14mm] (eng) at (10,0) {FPGA 엔진\\{\footnotesize 칩에 맞는 구조로}};
  \draw[<->] (rust) -- node[above,note] {대조} (ref);
  \draw[<->] (ref) -- node[above,note] {대조} (eng);
  \draw[<->] (rust.south) -- ++(0,-0.6) -| node[pos=0.25,below,note] {정지 목록을 순서까지 대조} (eng.south);
  \node[note,above=1mm of rust] {\code{crates/christopher-core}};
  \node[note,above=1mm of ref] {\code{rtl/src}};
  \node[note,above=1mm of eng] {\code{rtl/fpga}};
\end{tikzpicture}
\caption{세 층의 모델. 오른쪽으로 갈수록 실제 칩에 가깝고, 왼쪽으로 갈수록 읽고 믿기 쉽다.}
\label{fig:g-layers}
\end{figure}

기준 RTL(\code{rtl/src})은 처음부터 시뮬레이션용으로 썼다. 소프트웨어 모델을 가장 직접적으로 옮긴 것이어서
읽기 쉽고 Rust 코어와 대조하기 쉽다. 대신 표~\ref{tab:g-synth}의 항목을 여럿 갖고 있다. 나머지 연산
(\code{sum\_r \% 9'd26}), 변수로 676비트 가운데 하나를 고르는 접근, 선마다 통째로 놓인 스크램블러.
실제로 합성을 시도했을 때 이 코드는 합성 도구가 감당하지 못했다. 한 도구는 스크램블러 하나를 합성하다가
메모리를 13GB 넘게 쓰고 운영체제에 의해 종료되었고, 다른 칩을 대상으로 한 합성은 10분 안에 끝나지
않았다. 기준 RTL은 칩에 올리는 물건이 아니라 ``클럭 단위로 본 올바른 동작''의 정의다.

FPGA 엔진(\code{rtl/fpga})은 같은 동작을 칩에 맞는 구조로 다시 쓴 것이다. 엔진이 옳다는 근거는 엔진이
내놓는 정지 목록이 기준 RTL, 그리고 Rust 코어의 정지 목록과 \emph{순서까지} 같다는 것이다.

\section{검증 우선}

\Why 프로그램의 오류는 실행해 보고, 중간값을 찍어 보고, 고쳐서 다시 실행하면 된다. 몇 초면 한 바퀴 돈다.
회로는 그렇지 않다. 칩에 올리기까지 합성과 배치, 배선에 몇 분에서 몇 시간이 걸리고, 올린 뒤에는 칩 안의
신호를 들여다볼 수단이 마땅치 않다. 수십만 개의 신호 가운데 어느 것이 언제 틀렸는지 밖에서는 보이지 않는다.
그래서 회로는 칩에 올리기 \emph{전에}, 모든 신호를 볼 수 있는 시뮬레이션에서 최대한 확인한다.

확인하는 방법도 정해져 있다. 사람이 파형을 눈으로 보고 ``맞는 것 같다''고 하는 것은 검증이 아니다.
믿을 수 있는 기준 모델을 두고, 같은 입력을 회로와 모델에 넣어 출력을 기계적으로 비교한다. 비교는 자동이어야
하고, 입력은 사람이 떠올리지 못하는 경우까지 포함하도록 무작위로 대량 생성한다.

\General FIFO를 검증한다고 하자. 테스트벤치(회로를 시험하는 프로그램)는 클럭마다 무작위로 ``넣을까,
뺄까''를 정해 회로에 신호를 넣는다. 동시에 테스트벤치 안의 평범한 목록에도 같은 넣기와 빼기를 한다.
회로에서 값이 나올 때마다 목록에서 꺼낸 값과 비교한다. 수만 클럭을 돌려 한 번도 어긋나지 않으면 가득 찼을
때, 비었을 때, 넣기와 빼기가 같은 클럭에 일어날 때의 동작이 모두 맞는 것이다. 이런 경계 상황은 무작위
입력이 저절로 만들어 낸다.

\Ours 이 프로젝트의 테스트벤치는 Python으로 쓴다(\code{cocotb}라는 도구가 Python 코드와 시뮬레이터를
이어 준다). 기준 모델은 Rust 코어다. 흐름은 그림~\ref{fig:g-verify}와 같다.

\begin{figure}[ht]
\centering
\begin{tikzpicture}[x=1cm,y=1cm]
  \node[blkg,minimum width=28mm,minimum height=12mm] (core) at (0,1.2) {Rust 코어};
  \node[blko,minimum width=32mm,minimum height=12mm] (vec) at (4.4,1.2) {벡터 파일\\{\footnotesize 입력과 기대 출력}};
  \node[blkb,minimum width=30mm,minimum height=12mm] (tb) at (9.0,1.2) {테스트벤치\\{\footnotesize Python (cocotb)}};
  \node[blkn,minimum width=30mm,minimum height=12mm] (dut) at (9.0,-1.0) {회로\\{\footnotesize 시뮬레이터 안}};
  \node[blkr,minimum width=24mm,minimum height=12mm] (cmp) at (13.2,1.2) {비교\\{\footnotesize 다르면 실패}};
  \draw[->] (core) -- (vec); \draw[->] (vec) -- (tb); \draw[->] (tb) -- (cmp);
  \draw[->] ([xshift=-4mm]tb.south) -- node[left,note] {입력을 넣는다} ([xshift=-4mm]dut.north);
  \draw[->] ([xshift=4mm]dut.north) -- node[right,note] {출력을 읽는다} ([xshift=4mm]tb.south);
\end{tikzpicture}
\caption{검증의 흐름.}
\label{fig:g-verify}
\end{figure}

시험은 층마다 따로 있다. 아래층이 맞다는 것을 먼저 확인해야 위층이 실패했을 때 원인을 좁힐 수 있다.

\begin{itemize}
\item \code{scrambler}: 무작위로 고른 로터 순서와 위치 3,000건에서 26글자의 치환표를 비교한다.
\item \code{bombe\_core}: 벡터 파일의 각 위치에서 676비트 상태 \emph{전체}를 비교한다. 정지 여부만 비교하면
중간 상태가 틀려도 우연히 통과할 수 있기 때문이다.
\item \code{bombe\_sweep}: 로터 순서 하나를 전부 돌려 정지 목록을 순서까지 비교한다.
\end{itemize}

두 가지 요령이 더 있다. 테스트벤치는 회로가 쓰지 않아야 할 입력(비활성화된 메뉴 칸 등)에 일부러 무작위 값을
넣는다. 회로가 실수로 그 값에 의존하면 결과가 달라져 드러난다. 그리고 정지를 받아 가는 쪽의 수락 신호를
무작위로 늦춘다. 다음 절의 약속을 회로가 지키는지 보기 위해서다.

\section{핸드셰이크와 FIFO}\label{sec:handshake}

\Why 회로의 블록들은 속도가 다르다. 봄베는 드럼 위치 수천 곳을 지나는 동안 정지를 한 번 낼 수도, 연달아
낼 수도 있다. 정지를 받아 가는 쪽(PC와 통신하는 부분)은 훨씬 느리다. 보내는 쪽이 ``지금 값이 있다''고만
하고 다음으로 넘어가 버리면 받는 쪽이 놓칠 수 있다. 두 블록 사이에는 ``넘겼다, 받았다''를 확인하는 약속이
필요하다.

가장 널리 쓰이는 약속이 \term{valid/ready 핸드셰이크}{handshake}다. 신호는 둘이다. 보내는 쪽은 넘길 값이
있으면 \code{valid}를 올리고, 받는 쪽은 받을 수 있으면 \code{ready}를 올린다. 규칙은 하나다.
\emph{클럭 에지에서 \code{valid}와 \code{ready}가 둘 다 1이면 그 순간 값이 넘어간 것이다.} 보내는 쪽은
\code{valid}를 올린 뒤 값이 넘어갈 때까지 값을 바꾸지 않고 기다려야 한다. 받는 쪽이 준비되지 않아 보내는
쪽이 멈춰 서는 것을 \term{역압}{backpressure}이라 한다.

\General 그림~\ref{fig:g-handshake}는 값 두 개가 넘어가는 모습이다. 첫 번째 값 A는 받는 쪽이 준비되어 있어
곧바로 넘어간다. 두 번째 값 B는 받는 쪽이 두 클럭 동안 준비되지 않아 그동안 보내는 쪽이 값을 붙들고
기다린다.

\begin{figure}[ht]
\centering
\begin{tikztimingtable}[timing/slope=0.1,timing/coldist=4pt,xscale=2.4,yscale=1.25,semithick]
  \code{clk}   & 9{.5L .5H} \\
  \code{valid} & 1.5L 1H 1L 3H 2.5L \\
  \code{data}  & 1.5U 1D{A} 1U 3D{B} 2.5U \\
  \code{ready} & .5L 2H 3L 2H 1.5L \\
\extracode
  \begin{pgfonlayer}{background}
    \vertlines[cgray!50,dashed]{0.5,1.5,...,8.5}
    \fill[cgreenbg] (1.5,-7) rectangle (2.5,1.2);
    \fill[cgreenbg] (5.5,-7) rectangle (6.5,1.2);
    \fill[credbg] (3.5,-7) rectangle (5.5,1.2);
  \end{pgfonlayer}
  \node[note,text=cgreen] at (2.0,-7.9) {A 전달};
  \node[note,text=cred] at (4.3,-7.9) {기다림};
  \node[note,text=cgreen] at (6.3,-7.9) {B 전달};
\end{tikztimingtable}
\caption{valid/ready 핸드셰이크. 초록 구간의 끝 에지에서 값이 넘어간다. 빨간 구간에서는 보내는 쪽이 값을
유지한 채 기다린다.}
\label{fig:g-handshake}
\end{figure}

보내는 쪽이 멈춰 서는 것이 곤란하면 두 블록 사이에 FIFO를 둔다. FIFO는 고정 크기의 메모리와 카운터 두 개로
만든다(그림~\ref{fig:g-fifo}). 넣을 때는 쓰기 위치에 값을 쓰고 쓰기 위치를 하나 올리며, 뺄 때는 읽기 위치의
값을 내놓고 읽기 위치를 하나 올린다. 두 위치는 끝에 이르면 처음으로 돌아간다. 보내는 쪽은 FIFO가 가득 차지
않는 한 멈추지 않고, 받는 쪽은 자기 속도로 꺼내 간다.

\begin{figure}[ht]
\centering
\begin{tikzpicture}[x=1cm,y=1cm]
  \foreach \i/\v/\st in {0/{}/white,1/{}/white,2/C/cbluebg,3/D/cbluebg,4/E/cbluebg,5/{}/white,6/{}/white,7/{}/white} {
    \node[draw,fill=\st,minimum width=10mm,minimum height=9mm,font=\ttfamily\small] (c\i) at (\i*1.0,0) {\v};
    \node[font=\tiny,text=cgray] at (\i*1.0,-0.7) {\i};
  }
  \draw[->,cgreen,line width=1pt] (2,1.5) node[above,font=\footnotesize,text=cgreen] {읽기 위치} -- (c2.north);
  \draw[->,corange,line width=1pt] (5,1.5) node[above,font=\footnotesize,text=corange] {쓰기 위치} -- (c5.north);
  \draw[->] (-1.8,0) node[left,font=\small] {넣기} -- (c0.west);
  \draw[->] (c7.east) -- ++(1.2,0) node[right,font=\small] {빼기};
  \node[note,anchor=west,text width=36mm,align=left] at (9.9,0.2) {두 위치가 같으면 비어 있다.\\넣으면 쓰기 위치가, 빼면 읽기 위치가 하나 올라간다.};
\end{tikzpicture}
\caption{FIFO의 구조. 메모리 여덟 칸에 값 세 개(C, D, E)가 들어 있는 상태다.}
\label{fig:g-fifo}
\end{figure}

\Ours 기준 RTL의 \code{bombe\_sweep}은 정지를 valid/ready로 내보낸다. \code{stop\_valid}가 정지가 있음을
알리고, 받는 쪽이 \code{stop\_ready}로 수락할 때까지 탐색은 그 위치에 서 있다. 실제 봄베가 정지한 뒤
조작자가 값을 적고 다시 출발시킬 때까지 서 있던 것과 같은 동작이다. FPGA 엔진은 정지를 FIFO에 넣는다.
PC가 읽어 가는 것이 늦어도 FIFO가 차기 전까지는 탐색이 멈추지 않고, FIFO가 가득 차면 엔진이 멈춰 서서
정지를 잃지 않는다.

\section{정리}

\begin{itemize}
\item 시뮬레이션에서 옳은 코드가 모두 칩에 올릴 수 있는 코드는 아니다. 동작을 정의하는 기준 모델과 칩에
올리는 설계를 따로 두고 둘을 대조한다.
\item 회로는 칩에 올리기 전에 시뮬레이션에서 기준 모델과 자동으로, 대량으로 비교한다.
\item 속도가 다른 블록 사이에는 넘김을 확인하는 약속(valid/ready)을 두고, 필요하면 FIFO로 완충한다.
\end{itemize}
